\documentclass[11pt]{article}
\usepackage{amsmath,amssymb,amsthm,booktabs,array,geometry,hyperref}
\geometry{margin=1in}

\newtheorem{theorem}{Theorem}
\newtheorem{lemma}[theorem]{Lemma}
\newtheorem{proposition}[theorem]{Proposition}
\newtheorem{corollary}[theorem]{Corollary}
\newtheorem{definition}[theorem]{Definition}
\newtheorem{remark}[theorem]{Remark}

\title{Quantum Operation Costs Under the Closed 23-Operator Set:\\
Final Re-Audit with SuperBEST v5.2 Routing}
\author{Monogate Research}
\date{2026-04-21}

\begin{document}
\maketitle

\begin{abstract}
We re-compute exact F16-node counts for the full quantum computing catalog---including
time-evolution operators, von Neumann entropy, quantum relative entropy, Bures distance,
matrix square root, density-matrix logarithm, partition functions, and related quantities---using
the locked 23-operator SuperBEST set and v5.2 routing.
The 23-operator set (F16 plus EEA, EES, EEM, EED and LLS, LLA, LLD)
produces measurable reductions in five categories of quantum computation.
We document exact before-vs-after node counts and identify the structural reasons
for each collapse.
\end{abstract}

\section{Setup and Notation}

\subsection{The 23-Operator SuperBEST Set}

The closed operator taxonomy consists of:
\begin{itemize}
  \item \textbf{F16}: the 16-element orbit of $\mathrm{eml}(x,y)=e^x-\ln y$ under
    $G_{16}\cong(\mathbb{Z}/2)^4$, covering all sign/reciprocal variants of
    $\exp(\pm x)\,[\pm]\,\ln(\pm y)$ for $[\pm]\in\{+,-,\times,\div\}$.
  \item \textbf{EEA}: $e^x+e^y$ (1 node; was $4n$ via two exps + add)
  \item \textbf{EES}: $e^x-e^y$ (1 node)
  \item \textbf{EEM}: $e^{x+y}=e^x\cdot e^y$ (1 node)
  \item \textbf{EED}: $e^{x-y}=e^x/e^y$ (1 node; was $4n$)
  \item \textbf{LLS}: $\ln(x/y)=\ln x-\ln y$ (1 node; was $4n$)
  \item \textbf{LLA}: $\ln(xy)=\ln x+\ln y$ (1 node; was $4n$)
  \item \textbf{LLD}: $\log_y x=\ln x/\ln y$ (1 node; was $4n$)
\end{itemize}

\subsection{SuperBEST v5.2 Core Costs (Locked)}

\begin{center}
\begin{tabular}{lcc}
\toprule
Operator & Positive domain & General domain \\
\midrule
$e^x$ & $1n$ & $1n$ \\
$\ln x$ & $1n$ & $1n$ \\
$1/x$ & $1n$ & $1n$ \\
$x^n$ (pow) & $1n$ & $1n$ \\
$\sqrt{x}$ & $1n$ & $2n$ \\
$-x$ (neg) & $1n$ & $1n$ \\
$x\cdot y$ (mul) & $2n$ & $2n$ \\
$x-y$ (sub) & $2n$ & $2n$ \\
$x+y$ (add) & $2n$ & $2n$ \\
$x/y$ (div) & $2n$ & $2n$ \\
\midrule
\textbf{Positive total} & \textbf{15n} & --- \\
\textbf{General total} & --- & \textbf{22n} \\
\bottomrule
\end{tabular}
\end{center}

\subsection{Node-Count Measurement Convention}

An \emph{$n$-node expression} requires exactly $n$ applications of 23-operator primitives,
with inputs (eigenvalues $\lambda_i$, parameters $t$, $\beta$, etc.)\ counted as zero-cost leaves.
Shared sub-expressions are counted once (DAG semantics, not tree semantics) when explicitly stated.

\section{Quantum Matrix Operations}

\subsection{Matrix Square Root}

For a density matrix $\rho$ with spectral decomposition $\rho=\sum_i \lambda_i |i\rangle\langle i|$,
the matrix square root $\sqrt{\rho}$ acts spectrally: $\sqrt{\rho}=\sum_i\sqrt{\lambda_i}|i\rangle\langle i|$.

\begin{proposition}[Matrix sqrt cost]
For each eigenvalue $\lambda_i>0$, the scalar computation $\sqrt{\lambda_i}$ costs:
\begin{align*}
\text{v5.1 (16-op):} &\quad 2n \quad [\text{pow via two EML nodes}] \\
\text{v5.2 (23-op):} &\quad 1n \quad [\mathrm{EPL}(0.5,\lambda_i)=\exp(0.5\ln\lambda_i)=\sqrt{\lambda_i}]
\end{align*}
For a $d\times d$ density matrix: total savings $= d$ nodes.
\end{proposition}

\subsection{Density Matrix Logarithm}

$\ln\rho = \sum_i \ln(\lambda_i)|i\rangle\langle i|$. Each $\ln(\lambda_i)$ costs $1n$ in both versions.
\textbf{No change.}

\subsection{Matrix Exponential (Spectral Factor)}

The time-evolution operator $U(t)=e^{-iHt}$ acts on eigenvalues of $H$ via $e^{-\lambda_k t}$
(restricting to the real spectral domain of Hermitian $H$ acting on position eigenstates).

\begin{proposition}[Spectral decay factor]
For each real eigenvalue $\lambda_k$ and time parameter $t$:
\begin{align*}
\text{v5.1:} &\quad \mathrm{mul}(\lambda_k, t)\;[2n] + \mathrm{neg}\;[1n] + \exp\;[1n] = 4n \\
\text{v7:}   &\quad \mathrm{mul}(\lambda_k, t)\;[2n] + \mathrm{EED}(0,\cdot)\;[1n] = 3n
\end{align*}
where $\mathrm{EED}(0, z)=e^{0-z}=e^{-z}$.
\textbf{Savings: $1n$ per spectral factor.}
\end{proposition}

\begin{remark}
$\mathrm{EED}(0, z) = e^{-z}$ provides a direct one-node negated exponential gate.
This collapses the three-step \texttt{neg}+\texttt{exp}+\texttt{mul} pattern to two steps
whenever a negative-exponent evaluation is required.
\end{remark}

\section{Entropy and Information Quantities}

\subsection{Von Neumann Entropy}

$S(\rho)=-\mathrm{Tr}(\rho\ln\rho)=-\sum_i\lambda_i\ln\lambda_i$.

\begin{proposition}[Von Neumann entropy per term]
For each eigenvalue $\lambda_i$:
\begin{align*}
\text{v5.1 and v7:} &\quad \ln(\lambda_i)\;[1n] + \mathrm{neg}\;[1n] + \mathrm{mul}(\lambda_i,\cdot)\;[2n] = 4n
\end{align*}
\textbf{No change.} The form $-\lambda\ln\lambda$ is irreducible under F16 (no single 23-operator
covers the product of an identity gate with a log gate).
\end{proposition}

\subsection{Quantum Relative Entropy}

$D(\rho\|\sigma)=\mathrm{Tr}(\rho(\ln\rho-\ln\sigma))=\sum_i\lambda_i(\ln\lambda_i-\ln\mu_i)$,
where $\lambda_i,\mu_i$ are corresponding eigenvalues of $\rho,\sigma$.

\begin{theorem}[Quantum relative entropy collapse]
For each eigenvalue pair $(\lambda_i,\mu_i)$:
\begin{align*}
\text{v5.1:} &\quad \ln\lambda_i\;[1n] + \ln\mu_i\;[1n] + \mathrm{sub}\;[2n] + \mathrm{mul}(\lambda_i,\cdot)\;[2n] = 6n \\
\text{v7:}   &\quad \mathrm{LLS}(\lambda_i,\mu_i)\;[1n] + \mathrm{mul}(\lambda_i,\cdot)\;[2n] = 3n
\end{align*}
where $\mathrm{LLS}(\lambda_i,\mu_i)=\ln(\lambda_i/\mu_i)=\ln\lambda_i-\ln\mu_i$.
\textbf{Savings: $3n$ per eigenvalue pair; $3d$ total for a $d\times d$ system.}
\end{theorem}

\begin{proof}
LLS satisfies all four Primitive Gate Criteria (PGC): (1) outside F16 orbit, (2) depth-1 computable,
(3) saves $\geq 1n$ vs F16 routing, (4) canonically motivated as log-ratio.
The routing $\lambda_i\cdot\mathrm{LLS}(\lambda_i,\mu_i)$ is depth-2, taking eigenvalues as leaves.
\end{proof}

\subsection{Quantum Conditional Entropy}

$S(A|B)_\rho = S(\rho_{AB})-S(\rho_B) = -D(\rho_{AB}\|\mathbf{1}_A\otimes\rho_B)$.

Node count inherits from relative entropy and von Neumann entropy terms.
With LLS routing: reduces by $3n$ per joint eigenvalue pair over v5.1.

\subsection{Quantum Mutual Information}

$I(A:B)_\rho = S(\rho_A)+S(\rho_B)-S(\rho_{AB})$.
Each entropy computed at $4n$ per eigenvalue. \textbf{No new savings over v5.1} for
the entropy terms individually. Total savings come only from sub-expressions
that use sqrt or log-ratio patterns.

\section{Distance and Fidelity Measures}

\subsection{Bures Distance}

$d_B(\rho,\sigma)=\sqrt{2(1-F(\rho,\sigma))}$ where
$F(\rho,\sigma)=\left(\mathrm{Tr}\sqrt{\sqrt\rho\,\sigma\sqrt\rho}\right)^2$.

For commuting $\rho,\sigma$ (diagonal in the same basis):
$F(\rho,\sigma)=\left(\sum_i\sqrt{\lambda_i\mu_i}\right)^2$.

\begin{proposition}[Bures fidelity term]
For each pair $(\lambda_i,\mu_i)$:
\begin{align*}
\text{v5.1:} &\quad \mathrm{mul}(\lambda_i,\mu_i)\;[2n] + \sqrt{\cdot}\;[2n] = 4n \\
\text{v5.2:} &\quad \mathrm{mul}(\lambda_i,\mu_i)\;[2n] + \mathrm{EPL}(0.5,\cdot)\;[1n] = 3n
\end{align*}
\textbf{Savings: $1n$ per eigenvalue pair; $d$ total for $d\times d$ systems.}
\end{proposition}

\begin{remark}
An alternative routing: $\sqrt{\lambda_i\mu_i}=\sqrt{\lambda_i}\cdot\sqrt{\mu_i}$
using $\mathrm{EPL}(0.5,\lambda_i)\cdot\mathrm{EPL}(0.5,\mu_i)$ costs $1+1+2=4n$,
which is worse. The \texttt{mul}-first strategy is optimal.
\end{remark}

\subsection{Trace Distance}

$d_T(\rho,\sigma)=\tfrac12\|\rho-\sigma\|_1=\tfrac12\sum_i|\lambda_i-\mu_i|$.

For the scalar step $|\lambda_i-\mu_i|$: requires $\mathrm{sub}$ [2n] + absolute value
[$\notin$ ELC(ℝ) for general inputs, requires sign analysis].
No ELC routing improvement. \textbf{No change.}

\subsection{Quantum Jensen--Shannon Divergence}

$\mathrm{JSD}(\rho\|\sigma)=\tfrac12D\!\left(\rho\|\tfrac{\rho+\sigma}{2}\right)+\tfrac12D\!\left(\sigma\|\tfrac{\rho+\sigma}{2}\right)$.

Per eigenvalue tuple $(\lambda_i,\mu_i)$:
\[
m_i=\tfrac{\lambda_i+\mu_i}{2},\quad
\text{term}=\tfrac12\lambda_i\,\mathrm{LLS}(\lambda_i,m_i)+\tfrac12\mu_i\,\mathrm{LLS}(\mu_i,m_i).
\]
v7 count per tuple: $\mathrm{add}(\lambda_i,\mu_i)$ [2n] + scale by $\tfrac12$ [implicit constant mul, 0n with constant folding]
+ $\mathrm{LLS}(\lambda_i,m_i)$ [1n] + $\mathrm{mul}(\lambda_i,\cdot)$ [2n] + $\mathrm{LLS}(\mu_i,m_i)$ [1n]
+ $\mathrm{mul}(\mu_i,\cdot)$ [2n] + $\mathrm{add}$ [2n] = $10n$ per tuple.

v5.1 count: two LLS sub-expressions each expand to $1+1+2=4n$ instead of 1n, adding $6n$.
Total v5.1: $16n$ per tuple.
\textbf{Savings: $6n$ per eigenvalue tuple via two LLS collapses.}

\section{Partition Functions and Thermal States}

\subsection{Canonical Partition Function}

$Z(\beta)=\sum_{k=1}^n e^{-\beta E_k}$.

\begin{theorem}[Partition function EEA collapse]
For $n$ energy levels $E_1,\ldots,E_n$ and inverse temperature $\beta$:
\begin{align*}
\text{v5.1, per pair }(E_k,E_j): &\quad
  3n\;[\text{compute }-\beta E_k] + 3n\;[\text{compute }-\beta E_j]\\
  &\quad+ 1n\;[\exp(-\beta E_k)] + 1n\;[\exp(-\beta E_j)] + 2n\;[\text{add}] = 10n\\
\text{v7, per pair }(E_k,E_j): &\quad
  3n + 3n + 1n\;[\mathrm{EEA}(-\beta E_k,-\beta E_j)] = 7n
\end{align*}
\textbf{Savings: $3n$ per pair of energy levels; $3\lfloor n/2\rfloor$ total for $n$ levels.}
\end{theorem}

\begin{proof}
$\mathrm{EEA}(x,y)=e^x+e^y$ is a single node that subsumes the two \texttt{exp} calls and
the \texttt{add}, saving $3n$ per pair. The arguments $-\beta E_k$ are computed identically
in both versions via $\mathrm{EED}(0,\mathrm{mul}(\beta,E_k))$.
\end{proof}

\subsection{Free Energy}

$F(\beta)=-\beta^{-1}\ln Z(\beta)$.

Given $Z$: $\mathrm{recip}(\beta)$ [1n] + $\ln Z$ [1n] + $\mathrm{neg}$ [1n] + $\mathrm{mul}$ [2n] = $5n$.
No change from v5.1 at this level. Savings concentrated in $Z$ computation.

\subsection{Thermal State Entropy}

$S(\beta)=\ln Z+\beta\langle E\rangle_\beta$ where $\langle E\rangle_\beta=-\partial_\beta\ln Z$.

Node counts depend on representation of $\langle E\rangle_\beta$. For discrete spectrum:
$\langle E\rangle_\beta = Z^{-1}\sum_k E_k e^{-\beta E_k}$.
Each $E_k e^{-\beta E_k}$ costs: $3n$ [$e^{-\beta E_k}$ via EED] + $2n$ [$\mathrm{mul}(E_k,\cdot)$] = $5n$ per level.
\textbf{No change from v5.1 for individual terms.}
For paired sums using EEA on the exponential factors: see Section 5.1 approach.

\section{Lindblad and Open Quantum Systems}

\subsection{Lindblad Decay Factor}

Open quantum evolution produces factors of the form $e^{-\gamma_k t}(1-e^{-\gamma_j t})$
appearing in off-diagonal decoherence elements.

\begin{proposition}[Lindblad factor]
\begin{align*}
\text{v5.1:} &\quad 4n\;[e^{-\gamma_k t}] + 4n\;[e^{-\gamma_j t}] + 2n\;[\text{sub from 1}] + 2n\;[\text{mul}] = 12n \\
\text{v7:}   &\quad 3n\;[\mathrm{EED}(0,\gamma_k t)] + 3n\;[\mathrm{EED}(0,\gamma_j t)]\\
             &\quad+ 1n\;[\mathrm{EES}(0,\gamma_j t)=1-e^{-\gamma_j t}]\\
             &\quad+ 2n\;[\mathrm{mul}] = 9n
\end{align*}
\textbf{Savings: $3n$ per Lindblad decay pair.}
\end{proposition}

\begin{remark}
$\mathrm{EES}(0,z)=e^0-e^z=1-e^z$, so $\mathrm{EES}(0,-\gamma_j t)=1-e^{-\gamma_j t}$.
This collapses the three-node \texttt{neg}+\texttt{exp}+\texttt{sub-from-1} pattern.
\end{remark}

\subsection{Quantum Channel Kraus Operators}

Exponential factors in amplitude-damping channels: $\sqrt{1-e^{-\gamma t}}$ and $e^{-\gamma t/2}$.

\begin{align*}
e^{-\gamma t/2}: &\quad \mathrm{EED}(0,\mathrm{mul}(\gamma,t)/2)\;[2n+1n] = 3n \quad(\text{v7})\\
                 &\quad 4n\quad(\text{v5.1})\\
\sqrt{1-e^{-\gamma t}}: &\quad \mathrm{EES}(0,\gamma t)\;[3n]+\mathrm{EPL}(0.5,\cdot)\;[1n] = 4n\quad(\text{v7})\\
                        &\quad 3n\;[\text{compute }e^{-\gamma t}]+2n\;[\text{sub from 1}]+2n\;[\sqrt{\cdot}] = 7n\quad(\text{v5.1})
\end{align*}
\textbf{Savings: $1n$ for $e^{-\gamma t/2}$; $3n$ for $\sqrt{1-e^{-\gamma t}}$.}

\section{Complete Quantum Catalog: Before/After Table}

\begin{center}
\renewcommand{\arraystretch}{1.3}
\begin{tabular}{lcccc}
\toprule
Quantity & Domain & v5.1 (16-op) & v7 (23-op) & $\Delta$\\
\midrule
Matrix sqrt $\sqrt{\lambda}$ & per eigenvalue & $2n$ & $1n$ & $-1n$\\
Density log $\ln\lambda$ & per eigenvalue & $1n$ & $1n$ & $0$\\
Spectral decay $e^{-\lambda t}$ & per factor & $4n$ & $3n$ & $-1n$\\
vN entropy term $-\lambda\ln\lambda$ & per eigenvalue & $4n$ & $4n$ & $0$\\
Rel entropy term $\lambda\ln(\lambda/\mu)$ & per pair & $6n$ & $3n$ & $-3n$\\
Bures fidelity $\sqrt{\lambda\mu}$ & per pair & $4n$ & $3n$ & $-1n$\\
JSD term & per tuple & $16n$ & $10n$ & $-6n$\\
Partition function (per pair) & per 2 levels & $10n$ & $7n$ & $-3n$\\
Lindblad decay pair & per pair & $12n$ & $9n$ & $-3n$\\
Kraus $e^{-\gamma t/2}$ & per factor & $4n$ & $3n$ & $-1n$\\
Kraus $\sqrt{1-e^{-\gamma t}}$ & per factor & $7n$ & $4n$ & $-3n$\\
Thermal $\langle E\rangle_\beta$ term $E_k e^{-\beta E_k}$ & per level & $5n$ & $5n$ & $0$\\
Purity $\lambda^2$ & per eigenvalue & $1n$ & $1n$ & $0$\\
\midrule
\multicolumn{4}{l}{\textit{Outside ELC(ℝ): quantum gate unitaries ($e^{i\theta\sigma}$), Wigner function, etc.}} & $\infty$\\
\bottomrule
\end{tabular}
\end{center}

\section{Aggregate Savings for Standard Quantum Systems}

\subsection{Two-Qubit System ($d=4$)}

Representative computation: quantum relative entropy + Bures distance + matrix sqrt.
\begin{align*}
\text{v5.1:} &\quad 4\times 6n\;[\text{rel. ent.}] + 4\times 4n\;[\text{Bures}] + 4\times 2n\;[\text{sqrt}] = 24+16+8 = 48n\\
\text{v7:}   &\quad 4\times 3n + 4\times 3n + 4\times 1n = 12+12+4 = 28n\\
\text{Savings:} &\quad 20n \text{ for one relative-entropy + Bures + sqrt computation.}
\end{align*}

\subsection{$n$-Level System (Partition Function)}

For $n$ energy levels, pairwise EEA reduction:
\[
\Delta Z = 3\lfloor n/2\rfloor \text{ nodes saved.}
\]
For $n=100$ (atomic/molecular partition function): saves $150n$.

\subsection{Full Lindblad Trajectory ($d$ modes)}

Per time step: $d$ Lindblad decay pairs save $3d$ nodes;
$d$ Kraus $\sqrt{1-e^{-\gamma t}}$ terms save $3d$ nodes.
Total per time step: $6d$ nodes saved.
For $d=10$ modes over $T=1000$ steps: $60{,}000$ node-evaluations eliminated.

\section{Summary of New Collapses}

The five structural collapses produced by the 23-operator extension are:

\begin{enumerate}
  \item \textbf{Matrix sqrt}: $2n\to1n$ (sqrt correction in v5.2, EPL).
  \item \textbf{Quantum relative entropy}: $6n\to3n$ per pair (LLS collapses log-difference).
  \item \textbf{JSD}: $16n\to10n$ per tuple (two LLS collapses).
  \item \textbf{Partition function}: $10n\to7n$ per pair (EEA collapses exp-pair sum).
  \item \textbf{Lindblad / Kraus}: $12n\to9n$ and $7n\to4n$ (EED, EES collapse negated exponentials).
\end{enumerate}

Quantities unchanged: von Neumann entropy ($4n$ per eigenvalue), thermal energy
expectation ($5n$ per level), density log ($1n$), purity ($1n$). These are already
at structural lower bounds under the 23-operator taxonomy.

\end{document}
